\section{A curvature minor and a polynomial certificate for the null family}\label{app:certificate}

For the null branch of \eqref{eq:fields}, four real Killing spinors exist
throughout the family, and four exhaust the local solutions on an open dense
set. The two assertions use different certificates. A nonzero minor of the
supercurvature at one parameter value gives the generic upper bound. An
identity between the seed and deformed connections gives the lower bound,
including at parameters where the minor vanishes.

\paragraph{Frame and connection.}
Use the unsheared seed coframe $e^a$ and its dual $E_a$, with
$\eta_{ab}=\operatorname{diag}(-1,1,\ldots,1)$, $e^9=dx$, $e^{10}=dy$,
and orientation $e^0\wedge\cdots\wedge e^{10}$. Write $e^{abc}$ for an
ordered wedge. Define $C^{(0)}_{abc}=g_0([E_a,E_b],E_c)$: its only
nonzero entries, up to antisymmetry, are
\[
 C^{(0)}_{012}=-q_0,\qquad C^{(0)}_{345}=q_+,\qquad C^{(0)}_{678}=q_-,
 \qquad q_0^2=q_+^2+q_-^2,\quad q_i>0 .
\]
Thus $de^a=-\tfrac12 C^{(0)a}{}_{bc}\,e^b\wedge e^c$, where
$C^{(0)a}{}_{bc}=\eta^{ad}C^{(0)}_{bcd}$. Set $n=-E_0+E_1$,
$W=n+cE_{10}$, $w=W^\flat=e^0+e^1+c e^{10}$, $H_+=E_5$, $H_-=E_8$ and
$h=(1+\gamma^2c^2)^{1/2}$. Remove the constant homothety by taking
$(\hat g,\hat F)=(h^{-2/3}g,h^{-1}F)$. For $c\ne0$ the seed-orthogonal
projector onto $\operatorname{span}(W,H_+,H_-)$ is
\[
 \Pi^a{}_b=W^a w_b/c^2+H_+^a(e^5)_b+H_-^a(e^8)_b .
\]
The change of coframe, its inverse, and the corresponding frame vectors are
\begin{align*}
 M&=1+(h^{-1}-1)\Pi, & N&=M^{-1}=1+(h-1)\Pi,\\
 \vartheta^a&=M^a{}_b e^b, & e^b&=N^b{}_a\vartheta^a,
 &\widehat E_a&=N^b{}_a E_b .
\end{align*}
Indeed $\hat g=\eta_{ab}\vartheta^a\otimes\vartheta^b$: $M$ scales the
three selected directions by $h^{-1}$ and fixes their orthogonal complement.
To obtain the hatted flux components, make the substitution
$e^b=N^b{}_a\vartheta^a$ in the following seed-frame form:
\begin{align*}
 \hat F&=h^{-1}F_0+\gamma h^{-3}d(w\wedge e^5\wedge e^8)
               -\gamma h^{-1}c q_0 e^{3467},\\
 F_0&=(q_0e^{012}+q_+e^{345}+q_-e^{678})\wedge e^9 .
\end{align*}
In particular, the last term uses
$\iota_{H_-}\iota_{H_+}\iota_W\star_0 F_0=-c q_0e^{3467}$.

Define $C^d{}_{ab}$ by $[\widehat E_a,\widehat E_b]=C^d{}_{ab}\widehat E_d$
and put $C_{abc}=\eta_{cd}C^d{}_{ab}$. The new brackets and the spin
connection, with its two spin indices in the order used below, are
\begin{align*}
 C^d{}_{ab}&=M^d{}_k N^i{}_a N^j{}_b C^{(0)k}{}_{ij},\\
 2\omega_{abc}&=2\hat g(\widehat E_b,\nabla_{\widehat E_a}\widehat E_c)
             =C_{acb}-C_{cba}+C_{bac} .
\end{align*}
Use $\{\Gamma^a,\Gamma^b\}=2\eta^{ab}$, $\Gamma_a=\eta_{ab}\Gamma^b$, and
write $\operatorname{cl}(A)=A_{a_1\ldots a_p}\Gamma^{a_1\ldots a_p}/p!$
for the Clifford contraction in \eqref{eq:kse}. In this convention
the coefficient of each increasing wedge multiplies the corresponding
ordered gamma product without a factorial. The connection in \eqref{eq:kse}
and its curvature are $\mathcal D_a=\widehat E_a+T_a$ and
\begin{align*}
 T_a&=\tfrac12\sum_{b<c}\omega_{abc}\Gamma^b\Gamma^c
       +\tfrac1{12}\Gamma_a\operatorname{cl}(\hat F)
       -\tfrac14\operatorname{cl}(\iota_{\widehat E_a}\hat F),\\
 \mathcal R_{ab}&=[T_a,T_b]-C^d{}_{ab}T_d .
\end{align*}
The matrices are constant in the homogeneous frame; the components of a
parallel spinor need not be constant.

\paragraph{A nonzero minor.}
To specify rows and columns, complexify the Majorana module and index its
basis by $j=\sum_{k=0}^4 2^k j_k$, $0\le j\le31$. Bit $0$ is least
significant. On $|j\rangle$, $Z_k$ has eigenvalue $(-1)^{j_k}$, $X_k$
flips bit $k$, and $Y_k$ flips it with coefficient $i(-1)^{j_k}$.
Take the upper-index generators to be
\[
 \Gamma^0=i Z_0\cdots Z_4,\qquad
 \Gamma^{2k+1}=\Bigl(\prod_{l<k}Z_l\Bigr)X_k,\qquad
 \Gamma^{2k+2}=\Bigl(\prod_{l<k}Z_l\Bigr)Y_k .
\]
An increasing blade acts on a column rightmost generator first;
$\Gamma^0\cdots\Gamma^{10}=-1$. At
$(q_0,q_+,q_-,c,\gamma)=(10,6,8,3/4,1)$ one has $h=5/4$. Stack the
following rows in the order displayed and retain the columns $J$:
\[
 L=\begin{pmatrix}
       (\mathcal R_{01})_{0,\ldots,23\,;\,J}\\
       (\mathcal R_{02})_{0,2,4,6\,;\,J}
    \end{pmatrix},\qquad J=(0,1,\ldots,23,24,26,28,30).
\]
All denominators here are units modulo $257$. Reducing coefficients with
$i\mapsto16$ ($16^2\equiv-1$) gives $\det L\equiv140\pmod{257}$; the
leading $24\times24$ minor of $\mathcal R_{01}$ alone is $128\pmod{257}$.
Thus $L$ has rank $28$ in characteristic zero. Evaluation of a local
parallel spinor at one point is injective and takes values in the common
curvature kernel. Its real dimension is therefore at most four wherever
$\det L\ne0$. On the connected positive-radius cone times the real
$(c,\gamma)$ plane, $|\det L|^2$ is real analytic (see the regular frame
limit below). Its nonzero set is open and dense. A zero of this particular
minor is not a test for enhancement.

\paragraph{Transport of the seed subspace.}
In the same spinor components define
\begin{align*}
 J_0&=\Gamma^{012}(q_+\Gamma^{345}+q_-\Gamma^{678})/q_0,\\
 P_0&=\tfrac18(1+\Gamma^9)(1+\Gamma^{01})(1-J_0),\\
 K&=(\Gamma^0+\Gamma^1+c\Gamma^{10})\Gamma^5\Gamma^8,
 & S&=1+\frac{\gamma}{1+h}K .
\end{align*}
The three involutions in $P_0$ commute, and all nonscalar blades in the
expansion have zero trace. Thus $P_0$ is a real rank-four projector. To
see that its image carries four seed solutions, let $T_b^0$ denote $T_b$
at $\gamma=0$. Substituting $\Gamma^9P_0=P_0$ and
$(q_0\Gamma^{012}+q_+\Gamma^{345}+q_-\Gamma^{678})P_0=0$ into the
connection gives $T_b^0P_0=0$ in the AdS and flat directions. On sphere
$i$ it gives $T_b^0P_0=-(q_i/2)J_bP_0$, with $J_b=\Gamma_b\Gamma^{345}$
for $i=+$ and $J_b=\Gamma_b\Gamma^{678}$ for $i=-$. These $J_b$ commute
with $P_0$ and satisfy $[J_b,J_d]=-2J_f$ for cyclic sphere indices
$(b,d,f)$; the two sphere actions commute. Hence the restricted matrices
represent the seed brackets and the restricted connection is flat.

Since $K^2=-c^2$, $S$ is invertible even at $c=0$. The identity to be
certified in the hatted frame is
\begin{equation}
 \bigl(T_a S-S N^b{}_a T_b^0\bigr)P_0=0,
 \qquad a=0,\ldots,10 . \label{eq:certificate-transport}
\end{equation}
Indeed $\widehat E_a(S\epsilon_0)=S N^b{}_a E_b\epsilon_0$ in the two
homogeneous frames. Equation~\eqref{eq:certificate-transport} therefore
transports every local parallel seed section in $\operatorname{im}P_0$,
not merely its value at a point. This gives four independent solutions
without assuming that the common kernel of the curvature is parallel.

\paragraph{Exact verification over the family.}
For $c>0$ an AdS boost sending $n$ to $cn$ sends $W$ to $c(n+E_{10})$;
thus one can set $c=1$ and replace $\gamma$ by $c\gamma$. The boost
preserves $P_0$ (its spin generator is proportional to $\Gamma^{01}$).
Suppress the common positive frequency scale and set
\[
 \begin{gathered}
 (q_0,q_+,q_-)=(1+z^2,1-z^2,2z),\qquad 0<z<1,\\
 t=\frac{\gamma}{1+h},\qquad h=\frac{1+t^2}{1-t^2},\qquad
 \gamma=\frac{2t}{1-t^2},\qquad |t|<1 .
 \end{gathered}
\]
The omitted scale multiplies $C$, $\hat F$ and both connections by the
same constant and leaves $P_0$ unchanged. The entries of $M,N$ have
numerators of degree at most two over $1+t^2,1-t^2$, respectively. Put
$\Delta=(1+t^2)^3(1-t^2)^4$. The bracket uses one $M$ and two $N$'s;
the four-form uses four $N$'s multiplying $h^{-1}$, $\gamma h^{-3}$ or
$\gamma h^{-1}$. Thus $\Delta C$ and $\Delta\hat F$ have $t$-degree at
most $14$ and are linear in the frequencies. In the hatted frame $d$ of
an invariant form inserts one bracket and the Hodge star is constant.
The Ricci tensor is quadratic in the brackets by the connection formula
above. Consequently $\Delta^2$ clears the denominators in $d\hat F$,
$d\hat\star\hat F+\tfrac12\hat F\wedge\hat F$, and the Einstein residual
$\operatorname{Ric}_{ab}-\tfrac12\langle\iota_a\hat F,\iota_b\hat F\rangle
+\tfrac16\eta_{ab}|\hat F|^2$, leaving bidegree at most $(28,4)$ in
$(t,z)$. In \eqref{eq:certificate-transport}, $S=1+tK$ and $q_0P_0$ is
linear in the frequencies; multiplication by $\Delta q_0$ gives bidegree
at most $(15,4)$.

On the Cartesian grid
\[
 t=j/17\quad(-14\le j\le14),\qquad z=m/7\quad(1\le m\le5),
 \qquad j,m\text{ integers},
\]
all bosonic residuals and all eleven projected differences vanish in
exact arithmetic. For the latter test one may multiply rational Clifford
blades by reordering the indices with $\{\Gamma^a,\Gamma^b\}=2\eta^{ab}$,
replacing blades of degree greater than five by their complements using
$\Gamma^0\cdots\Gamma^{10}=-1$. The remaining blades form a basis of the
irreducible matrix algebra. No modular rank test enters this step.
For each fixed $z$ there are $29$ distinct roots in $t$; each resulting
coefficient has five distinct roots in $z$. The degree bounds therefore
prove polynomial identities. The extension to $c\le0$ follows by real
analyticity, since the only apparent frame poles satisfy
\[
 \frac{h-1}{c^2}=\frac{\gamma^2}{1+h},\qquad
 \frac{h^{-1}-1}{c^2}=-\frac{\gamma^2}{h(1+h)} .
\]
This certifies the lower bound at every regular parameter value, but does
not determine the exceptional set for the upper bound. In particular the
undeformed seed at $\gamma=0$ has sixteen spinors. The accompanying supplementary source includes the exact-arithmetic implementation. Both the minor and the interpolation can be reproduced with
\verb|python3 -B check_null.py --family| from the directory containing
the script (Python 3.10 or later, no external packages).
